% !TEX root = ../main.tex
\chapter{値・型・関数}
\label{chap:ch02_values_types_functions}

\begin{chapterabstract}
本章では、Lean の基本的な型と値、そして \texttt{def} による関数定義を扱う。
前章で読み方を固定した \texttt{:}、\texttt{:=}、\texttt{\#check}、\texttt{\#eval} を使いながら、
\texttt{Nat}、\texttt{Bool}、\texttt{String} の値を評価し、引数を取る関数を定義し、その型を確認する。
命題や証明にはまだ進まない。
\end{chapterabstract}

\begin{learninggoals}
\item \texttt{Nat}、\texttt{Bool}、\texttt{String} の値の書き方と評価結果を読める。
\item \texttt{def} で、引数のない値と、引数を取る関数を定義できる。
\item 関数の型 \texttt{A -> B} を「\texttt{A} を受け取り \texttt{B} を返す」と読める。
\end{learninggoals}

\section{この章を読む理由}

ソフトウェアでは、値の種類を取り違えると、動いていても意味が崩れる。
数値、真偽値、文字列は、それぞれ別の型として扱われる。
Lean では、この区別が型として明示され、関数にも入力と出力の型が付く。
本章では、その基本型と関数定義を、確認コマンドで結果を見ながら身につける。

図\ref{fig:ch02-types-values-functions}では、型が値を分類し、関数が入力型の値を出力型の値へ移す関係を示す。

\FigurePlaceholder{fig-ch02-types-values-functions.png}{0.90}{型・値・関数}{fig:ch02-types-values-functions}

\section{基本型の値}

\texttt{Nat} は自然数、\texttt{Bool} は真偽値、\texttt{String} は文字列の型である。
それぞれの値は、\texttt{0}、\texttt{1} のような数値、\texttt{true}、\texttt{false}、そして二重引用符で囲んだ文字列として書く。
演算子は型ごとに異なり、\texttt{+} は加算、\texttt{\&\&} は論理積、\texttt{++} は文字列連結である。

\begin{listing}[htbp]
\begin{minted}{lean4}
#eval (3 : Nat) + 4
#eval true && false
#eval "lean" ++ "4"
\end{minted}
\caption{\texttt{Nat}・\texttt{Bool}・\texttt{String}}
\label{lst:ch02_values}
\end{listing}

コード\ref{lst:ch02_values}では、\texttt{(3 : Nat) + 4} が \texttt{7} に、\texttt{true \&\& false} が \texttt{false} に、\texttt{"lean" ++ "4"} が \texttt{"lean4"} に評価される。
型が違えば使える演算子も違う、という点が読み取れれば十分である。

\section{\texttt{def} による値と関数}

\texttt{def} は、名前付きの定義を作る構文である。
引数のない値も、引数を取る関数も、同じ \texttt{def} で定義できる。
形は \texttt{def 名前 (引数 : 型) : 返り値の型 := 本体} である。

\begin{listing}[htbp]
\begin{minted}{lean4}
def answer : Nat :=
  42

def addOne (n : Nat) : Nat :=
  n + 1

def add (m : Nat) (n : Nat) : Nat :=
  m + n

#eval answer
#eval addOne 5
#eval add 2 3
\end{minted}
\caption{\texttt{answer}・\texttt{addOne}・\texttt{add}}
\label{lst:ch02_def}
\end{listing}

コード\ref{lst:ch02_def}では、\texttt{answer} は \texttt{Nat} 型の値、\texttt{addOne} は引数を一つ取る関数、\texttt{add} は引数を二つ取る関数である。
複数の引数は \texttt{(m : Nat) (n : Nat)} のように並べて書き、呼び出すときも \texttt{add 2 3} と空白で並べる。

\begin{warningbox}
\texttt{def} は可変変数への代入ではない。
\texttt{def answer : Nat := 42} は「\texttt{answer} という名前を \texttt{42} に結びつける」という意味であり、
あとから別の値に更新する、という読み方はしない。
\end{warningbox}

\section{関数の型}

関数にも型がある。
\texttt{A -> B} は「\texttt{A} 型の値を受け取り、\texttt{B} 型の値を返す関数」の型である。
前節の \texttt{addOne} は \texttt{Nat} を受け取り \texttt{Nat} を返すので、その型は \texttt{Nat -> Nat} である。
\texttt{\#check} で関数の型を確認できる。

\begin{listing}[htbp]
\begin{minted}{lean4}
def addOne (n : Nat) : Nat :=
  n + 1

def add (m : Nat) (n : Nat) : Nat :=
  m + n

#check addOne
#check add
\end{minted}
\caption{\texttt{addOne} と \texttt{add} の型}
\label{lst:ch02_function_type}
\end{listing}

コード\ref{lst:ch02_function_type}では、\texttt{\#check addOne} が \texttt{addOne} の型を、\texttt{\#check add} が \texttt{add} の型を表示する。
\texttt{addOne} は \texttt{Nat -> Nat}、\texttt{add} は二つの \texttt{Nat} を受け取る関数として読める。

\begin{leanbox}
\texttt{A -> B -> C} という型は、右に結合して \texttt{A -> (B -> C)} と読む。
つまり \texttt{add} は「\texttt{Nat} を受け取り、\texttt{Nat -> Nat} を返す関数」とも読める。
本章では「複数の引数を順に受け取る」と読めれば十分である。
\end{leanbox}

\section{本章のまとめ}

\begin{summarybox}
要点と記法
\begin{itemize}
\item \texttt{Nat} / \texttt{Bool} / \texttt{String} \quad 自然数・真偽値・文字列の型。
\item \texttt{+} / \texttt{\&\&} / \texttt{++} \quad 加算・論理積・文字列連結。
\item \texttt{def 名前 (引数 : 型) : 返り値の型 := 本体} \quad 関数定義の基本形。
\item \texttt{A -> B} \quad 「\texttt{A} を受け取り \texttt{B} を返す」関数の型。
\end{itemize}
\end{summarybox}

\section{次章への接続}

ここまでで、値・型・関数を読み書きできるようになった。
次章では、これらの関数について「\texttt{hasDiscount true} の結果は \texttt{true} と等しい」のような\emph{等式}を扱う。
そこで初めて、\texttt{=} を「代入」ではなく「主張（命題）」として読む練習を始める。
